% !TEX root = userguide.tex

\def\headdate{17th November 2022}

\section{The Poisson equation}

See Sec.~\ref{sec:poissonweak} for the theory.

\subsection{A Poisson problem on a unit square with Dirichlet
            and Neumann boundary conditions. Axisymmetric}

\label{sec:poisson13}
We consider the Poisson problem on a unit square, like in
Sec.~\ref{sec:poisson2}, however we assume the solution is axisymmetrical.
The $(x,y)$ coordinate changes now to $(z,r)$. The curve \texttt{C1} in
Fig.~\ref{fig:unitsquare} is at $r=R_0$ and \texttt{C3} at $r=R_1$. We assume
the flux is given at $r=R_0$ (curve \texttt{C1}):
\begin{equation}
   -\pderiv{u}{n} = h_\text{0}\text{ on \texttt{C1}}
\end{equation}
The following solution fulfills the Poisson equation:
\begin{equation}
   u(r,z) = u_1 - h_0R_0\ln\frac{R_1}r
\end{equation}
where $u_1$ is the constant value of $u$ at $r=R_1$.

The source is in the file \texttt{poisson13.f90}.

\subsection{A Poisson problem on a unit square with Dirichlet
            and Robin boundary conditions}

\label{sec:poisson12}
We consider the Poisson problem on a unit square, like in
Sec.~\ref{sec:poisson2}. 
We change the boundary condition on curve \texttt{C2} to a Robin boundary
condition, meaning that we have to specify (see Sec.~\ref{sec:bounc})
\begin{equation}
   -\pderiv{u}{n} - \sigma u = h_\text{R}\text{ on \texttt{C2}}
\end{equation}
We use $\sigma=1$ in the example.
We substitute the following exact solution for $u$:
\begin{equation}
     u=\cos(\pi x)\cos(\pi y) + x^3y^3
\end{equation}
and find the corresponding $f$:
\begin{equation}
     f=2\pi^2\cos(\pi x)\cos(\pi y) - 6(xy^3+x^3y)
\end{equation}
We specify $h_R$ by a function \texttt{func}. The exact value on \texttt{C2}
is given by
\begin{equation}
   h_R= \pi\sin(\pi x)\cos(\pi y) - 3x^2y^3 - \cos(\pi x)\cos(\pi y) - x^3y^3
\end{equation}
The source is in the file \texttt{poisson12.f90} and the difference with respect
to \texttt{poisson2.f90} will be discussed. The changes/added lines are:
\bigskip
\begin{listing}{22}
    bfuncnr=13            ! function number for the flux h_R on the Robin b
\end{listing}
\bigskip
\begin{listing}{26}
    sigma = 1._dp     ! coefficient sigma in the Robin bc
\end{listing}
\bigskip
\begin{listing}{48}
    coefficients%i(15) = bfuncnr
\end{listing}
\bigskip
\begin{listing}{51}
`\end{listing}
\bigskip
\begin{listing}{99}
  call add_boundary_elements ( mesh, problem, rhsd, curve=2, &
    sysmatrix=sysmatrix, elemsub=poisson_natboun, coefficients=coefficients )

\end{listing}
\begin{description}
   \item[line 22] changed \texttt{vfuncnr} to \texttt{bfuncnr}.
   \item[line 26] define $\sigma$
   \item[line 48] put the function number for $h_R$ into
\texttt{coefficients\%i}.
   \item[line 51] put $\sigma$ into \texttt{coefficients\%r}.
   \item[line 99 and 100]here the natural boundary integrals are
 added using the element subroutine
\texttt{poisson\_natboun}
and using the function \texttt{func}. The function number
is supplied via the coefficients. Both a system matrix and vector are build,
by specifying \texttt{sysmatrix=sysmatrix}.
\end{description}

\subsection{A Poisson problem on two unit squares with Dirichlet boundary 
            conditions. Coupling of the two domains using a weak
            constraint on an object.}
\label{sec:poisson8}

We consider the Poisson problem on two unit squares. The problem is similar to
the problem in Sec.~\ref{sec:poisson1}, but now extended to two squares and
shifted a distance of 0.2 to the right. The two separate squares are coupled
using a weak coupling, implemented using Lagrange multipliers one order lower
than the primary variable (see \cite{SeSu00} for the theory).

The source is in the file \texttt{poisson8.f90}. 

In Fig.~\ref{fig:mesh_mortar} the meshes used are shown, and typical results
are shown in Fig.~~\ref{fig:sol_mortar}.
\begin{figure}[phb!]
   \begin{center}
   \includegraphics[scale=0.5]{mesh_left}\medskip\\
   \includegraphics[scale=0.5]{mesh_right}
   \end{center}
   \caption{Mesh with the Lagrange multipliers defined on the left part
(top) and the right part of the mesh (bottom), respectively. The red dots are
the nodes of the quadratic line elements of the object
at the interface. The Lagrange
multipliers are only defined at the end nodes of an element (continuous linear
interpolation).}
    \label{fig:mesh_mortar}
\end{figure}
\begin{figure}[pthb!]
   \begin{center}
   \includegraphics[scale=0.5]{sol_left}\medskip\\
   \includegraphics[scale=0.5]{sol_right}
   \end{center}
   \caption{Results with the Lagrange multipliers defined on the left part
(top) and the right part of the mesh (bottom), respectively.}
    \label{fig:sol_mortar}
\end{figure}

\subsection{A Poisson problem on a unit cube with Dirichlet
            and Neumann boundary conditions}

\label{sec:poisson5}

We consider the Poisson problem on a unit cube. The problem is similar to
the problem in Sec.~\ref{sec:poisson2}, but now extended to 3D. The Neumann
boundary condition, where that we have to specify $\plderiv{u}{n}$,
is now on surface 3 (see Fig.~\ref{fig:curves3D}).
\begin{figure}[ht!]
   \begin{center}
   \includegraphics[scale=0.5]{curves}
   \end{center}
   \caption{Plot of points, curves and surfaces on the unit cube.}
    \label{fig:curves3D}
\end{figure}
We substitute the following exact solution for $u$:
\begin{equation}
     u=\cos(\pi x)\cos(\pi y)\cos(\pi z) + x^3y^3z^3
\end{equation}
and find the corresponding $f$:
\begin{equation}
     f=3\pi^2\cos(\pi x)\cos(\pi y)\cos(\pi z) - 6(xy^3z^3+x^3yz^3+x^3y^3z)
\end{equation}
Specifying $\plderiv{u}{n}$ is easy for the unit cube, but for a
curvilinear domain it is very difficult. In that case, it is much easier to
specify $\nabla u$ by a function
and let the boundary element compute $\plderiv{u}{n}=\vec n\cdot\nabla u$.
This is what we do here as well and specify $\nabla u$ by a vector function
\texttt{vfunc}. The exact $\nabla u$ is given by
\begin{equation}
   \nabla u=
\begin{pmatrix}
     -\pi\sin(\pi x)\cos(\pi y)\cos(\pi z) + 3x^2y^3z^3 \\
     -\pi\cos(\pi x)\sin(\pi y)\cos(\pi z) + 3x^3y^2z^3 \\
     -\pi\cos(\pi x)\cos(\pi y)\sin(\pi z) + 3x^3y^3z^2
\end{pmatrix}
\end{equation}
The source is in the file \texttt{poisson5.f90}. The mesh is generated using
the routine \texttt{hexahedron} from the \tfem\ mesh generator.
\begin{figure}[pt!]
   \begin{center}
   \includegraphics[scale=0.5]{mesh3D}
   \end{center}
   \caption{Mesh on a unit cube using $5\times5\times5$ elements.}
    \label{fig:mesh3D}
\end{figure}

Typing at the command line 
\begin{quote}
    \texttt{poisson5}
\end{quote}
will show something like
\begin{quote}
  \texttt{1.520878241239032E-04}
\end{quote}
on your screen. In Fig.~\ref{fig:datasurface3} the solution on surface 3 is
plotted.
\begin{figure}[pt!]
   \begin{center}
   \includegraphics[scale=0.5]{color}
   \end{center}
   \caption{Plot of data on surface 3 for \texttt{poisson5}.}
    \label{fig:datasurface3}
\end{figure}

The examples \texttt{poisson27.f90} and \texttt{poisson28.f90} are similar to \texttt{poisson5.f90}, but now the elements are prisms with linear and quadratic base functions, respectively. The meshes are generated using gmsh with the files \texttt{cube1.geo} and \texttt{cube2.geo} as input.

The examples \texttt{poisson29.f90} and \texttt{poisson30.f90} are similar to
\texttt{poisson5.f90}, but now hybrid meshes are used where both tetrahedra and
pyramids are used. The example  \texttt{poisson29.f90} uses linear base
functions, whereas  \texttt{poisson30.f90} uses quadratic ones. The meshes are generated using gmsh with the files \texttt{cube3.geo} and \texttt{cube4.geo} as input.

The example \texttt{poisson31.f90} is similar to \texttt{poisson30.f90}, but
now uses a mapped integration rule on the pyramid elements.

\subsection{A Poisson problem on a unit cube with Dirichlet
            and periodic boundary conditions; collocation}
\label{sec:poisson6}

We consider the Poisson problem on a unit cube. The problem is similar to
the problem in Sec.~\ref{sec:poisson5}, but now with periodical boundary
conditions in $x$-direction using collocation of the constraint.

We substitute the following exact solution for $u$:
\begin{equation}
     u=\cos(2\pi x+1)\cos(\pi y)\cos(\pi z) + y^3z^3
\end{equation}
and find the corresponding $f$:
\begin{equation}
     f=6\pi^2\cos(2\pi x+1)\cos(\pi y)\cos(\pi z) - 6(yz^3+y^3z)
\end{equation}
The source is in the file \texttt{poisson6.f90}.
Typing at the command line 
\begin{quote}
    \texttt{poisson6}
\end{quote}
will show something like
\begin{quote}
   \texttt{8.762962829161891E-04}
\end{quote}
on your screen.

\subsection{A Poisson problem on a unit cube with Dirichlet
            and periodic boundary conditions; weak connection}

We consider the Poisson problem on a unit cube. The problem is similar to
the problem in Sec.~\ref{sec:poisson6}, but now 
using a weak form of the constraint.

The source is in the file \texttt{poisson7.f90}.
Typing at the command line 
\begin{quote}
    \texttt{poisson7}
\end{quote}
will show something like
\begin{quote}
   \texttt{2.958958242691212E-03}
\end{quote}
on your screen.


\subsection{A Poisson problem on a unit square with Dirichlet boundary 
            conditions. Line source. Build system on an object.}
\label{sec:poisson9}

We consider the Poisson problem on a unit square. On the outside boundary we
impose $u=0$. We assume the following singular source term on a line $C$
 inside the domain:
\[
   f=(a+bu)\,\delta(n)
\]
where $\delta(n)$ is a delta function in the normal direction of the line.
A physical meaning is, for example, a temperature dependent heat source
concentrated on a line (a line source).
The right-hand side of \eqref{eq:weakformpoisson} now becomes:
\[
(v,f)=\int_C v(a+bu)\D s=\int_C va\D s+\int_C bvu\D s
\]
where $s$ is a curve coordinate along the line $C$. The two terms lead to 
contributions to the system vector and system matrix, respectively. The
integrals are computed on an object consisting of line elements. The object is
constructed from a mesh consisting of line elements.

The source is in the file \texttt{poisson9.f90}. 

In Fig.~\ref{fig:poisson_linesource} a typical result is shown.
\begin{figure}[phbt!]
   \begin{center}
   \includegraphics[scale=0.5]{poisson_linesource}
   \end{center}
   \caption{Solution of the poisson equation with a line source ($a=1$, $b=-1$)
on a unit square.}
    \label{fig:poisson_linesource}
\end{figure}

\subsection{A Poisson problem on a unit square with Dirichlet boundary 
            conditions using spectral or high-order elements}
\label{sec:poisson20}

The examples are similar to \texttt{poisson1.f90} (see
Sec.~\ref{sec:poisson1}), but now with spectral or high-order elements\footnote{In \tfem\ high-order refers to elements with high-order polynomial interpolation using standard Gauss integration and regular (equidistant) nodal point distribution. Spectral elements also use high-order polynomial interpolation, but Gauss-Lobatto points for both integration and nodal point distribution.}
The source files are \texttt{poisson20.f90} (spectral) and \texttt{poisson34.f90} (high-order quadrilaterals).
The files and can be obtained by typing at the
command line:
\begin{quote}
   \texttt{getexample poisson}
\end{quote}

\subsection{A Poisson problem on a curved quadrilateral region with Dirichlet 
            and Neumann boundary conditions using spectral or high-order elements} 
\label{sec:poisson21}

This examples are similar to \texttt{poisson3.f90} (see
Sec.~\ref{sec:poisson3}), but now with spectral or high-order elements.
The source files are \texttt{poisson21.f90} (spectral) and \texttt{poisson35.f90} (high-order quadrilaterals).
The files and can be obtained by typing at the
command line:
\begin{quote}
   \texttt{getexample poisson}
\end{quote}

\subsection{A Poisson problem on a curved quadrilateral region with Dirichlet 
            and Neumann boundary conditions. Derived quantities. Spectral or high-order
            elements (quadrilaterals and triangles)} 
\label{sec:poisson22}

This examples are similar to \texttt{poisson4.f90} (see
Sec.~\ref{sec:poisson4}), but now with spectral or high-order elements (quadrilaterals and triangles).
The source files are \texttt{poisson22.f90} (spectral), \texttt{poisson36.f90} (high-order quadrilaterals) and \texttt{poisson37.f90} (high-order triangles).
The files and can be obtained by typing at the
command line:
\begin{quote}
   \texttt{getexample poisson}
\end{quote}

\subsection{A Poisson problem on a unit cube with Dirichlet
            and periodic boundary conditions; collocation; non-symmetric solver MA41 with \metis\ renumbering}
\label{sec:poisson24}

We consider the Poisson problem on a unit cube. The problem is similar to
the problem in Sec.~\ref{sec:poisson6}, but now using the non-symmetric 
solver MA41 (see Sec.~\ref{sec:hslsparse}) together with \metis\ renumbering (see Sec.~\ref{sec:metis5}). The purpose of this example is to illustrate how to use \metis\ with MA41, which can speedup the solver significantly ($3.6$ times for a $30\times30\times30$ mesh!). The source is in the file \texttt{poisson24.F90}
and can be obtained by typing at the
command line:
\begin{quote}
   \texttt{getexample poisson}
\end{quote}

\subsection{A Poisson problem on a unit cube with Dirichlet
            and periodic boundary conditions; collocation; Sparskit2 with \metis\ renumbering of the degrees of freedom}
\label{sec:poisson25}

We consider the Poisson problem on a unit cube. The problem is similar to
the problem in Sec.~\ref{sec:poisson6}, but now using the Sparskit2/sk\_solve (see Sec.~\ref{sec:sksolve})
together with \metis\ renumbering of the degrees of freedom and/or nodes (see Sec.~\ref{sec:metis5}). The source is in the file \texttt{poisson25.F90}
and can be obtained by typing at the
command line:
\begin{quote}
   \texttt{getexample poisson}
\end{quote}

\metis\ renumbering of the degrees of freedom instead of only nodes renumbering can positively influence 
the memory/CPU requirements for building the preconditioner and/or the convergence rate. However, it is not guaranteed. As always with iterative solvers: experiment with the available parameters.

The example in \texttt{poisson26.F90} is similar to \texttt{poisson25.F90}, however the preconditioner has
been changed from \texttt{ilut} to \texttt{ilutp}, which uses column pivoting. The effect is rather small, at least for this example.
